\documentclass[10pt]{article}
\usepackage[ngerman]{babel}
\usepackage[utf8]{inputenc}
\usepackage[T1]{fontenc}
\usepackage{amsmath}
\usepackage{amsfonts}
\usepackage{amssymb}
\usepackage[version=4]{mhchem}
\usepackage{stmaryrd}

\title{Hauptprüfung Abiturprüfung 2024 - Leistungsfach }

\author{Alexander Schwarz\\
www.mathe-aufgaben.com}
\date{}


\begin{document}
\maketitle
 Baden-Württemberg

\section*{Aufgabe III 2}
Eine Berghütte wird von Gästen mit Mountainbike und ohne Mountainbike besucht. Es gibt Gäste, die eigene Verpflegung mitbringen. Erfahrungsgemäß kommen 44\% der Gäste mit dem Mountainbike. 73\% der Gäste bringen keine eigene Verpflegung mit. $11 \%$ der Gäste kommen mit dem Mountainbike und bringen eigene Verpflegung mit.\\
Ein Gast wird zufällig ausgewählt. Betrachtet werden die folgenden Ereignisse: M: „Der Gast kommt mit dem Mountainbike."\\
V : „Der Gast bringt eigene Verpflegung mit."\\
a) Stellen Sie den Sachverhalt in einer vollständig ausgefüllten Vierfeldertafel dar.\\
b) Untersuchen Sie, ob der Anteil der Gäste, die eigene Verpflegung mitbringen, unter den Gästen mit Mountainbike gleich groß ist wie unter den Gästen ohne Mountainbike.\\
c) Berechnen Sie die Wahrscheinlichkeit des Ereignisses $\overline{\mathrm{M}} \cup \mathrm{V}$.\\
d) Geben Sie jeweils einen Wert für die Parameter $a$ und $b$ an ( $a>0, b>0$ ), sodass mit dem Term $a \cdot 0,44 \cdot 0,56^{b}+0,56^{10}$ die Wahrscheinlichkeit eines Ereignisses im Sachzusammenhang berechnet werden kann. Beschreiben Sie das zugehörige Zufallsexperiment und das Ereignis.\\
(4 BE)\\
Auf der Berghütte werden Bananen, Birnen und Äpfel angeboten. Die Masse der Bananen wird als normalverteilt mit einem Erwartungswert von 120 g und einer Standardabweichung von 10 g angenommen.\\
e) Berechnen Sie die Wahrscheinlichkeit dafür, dass die Masse einer zufällig ausgewählten Banane um höchstens 10\% vom Erwartungswert abweicht.\\
f) Die Wahrscheinlichkeit, dass eine zufällig ausgewählte Banane eine größere Masse als $m$ hat, beträgt 69\% (m in Gramm). Ermitteln Sie den Wert von m.\\
g) Die Masse der Birnen wird ebenfalls als normalverteilt mit einem ganzzahligen Erwartungswert in Gramm und einer Standardabweichung von 25 g angenommen. Die Wahrscheinlichkeit dafür, dass eine zufällig ausgewählte Birne eine Masse von höchstens 121 g hat, beträgt 10\%. Bestimmen Sie die Wahrscheinlichkeit dafür, dass eine zufällig ausgewählte Birne mindestens 180 g wiegt.\\
h) Die Masse eines Apfels wird als normalverteilt mit $\mu_{\mathrm{A}}$ und $\sigma_{\mathrm{A}}$ angenommen. Berechnen Sie die Wahrscheinlichkeit dafür, dass von vier zufällig ausgewählten Äpfeln zwei eine kleinere Masse als $\mu_{\mathrm{A}}$ und zwei eine größere Masse als $\mu_{\mathrm{A}}$ haben.\\
(3 BE)

\section*{Lösungen}
\section*{Aufgabe III 2}
a) Aus dem Text ergibt sich:\\
$\mathrm{P}(\overline{\mathrm{V}})=0,73$ und $\mathrm{P}(\mathrm{M})=0,44$ und $\mathrm{P}(\mathrm{M} \cap \mathrm{V})=0,11$

\begin{center}
\begin{tabular}{|c|c|c|c|}
\hline
 & V & $\overline{\mathrm{V}}$ &  \\
\hline
M & 0,11 & 0,33 & 0,44 \\
\hline
$\overline{\mathrm{M}}$ & 0,16 & 0,4 & 0,56 \\
\hline
 & 0,27 & 0,73 & 1 \\
\hline
\end{tabular}
\end{center}

b) Anteil der Gäste, die eigene Verpflegung mitbringen, unter den Gästen mit

Mountainbike: $\mathrm{P}_{\mathrm{M}}(\mathrm{V})=\frac{\mathrm{P}(\mathrm{M} \cap \mathrm{V})}{\mathrm{P}(\mathrm{M})}=\frac{0,11}{0,44}=\frac{1}{4}$\\
Anteil der Gäste, die eigene Verpflegung mitbringen, unter den Gästen ohne\\
Mountainbike: $P_{\bar{M}}(\mathrm{~V})=\frac{\mathrm{P}(\overline{\mathrm{M}} \cap \mathrm{V})}{\mathrm{P}(\overline{\mathrm{M}})}=\frac{0,16}{0,56}=\frac{2}{7}$\\
Die Anteile sind nicht gleich.\\
c) $P(\bar{M} \cup V)=0,11+0,16+0,4=0,67$\\
d) Der Term a $\cdot 0,44 \cdot 0,56^{b}+0,56^{10}$ stellt die Summe zweier Formeln von Bernoulli dar.

Aus $0,56^{10}$ ergibt sich $n=10$ Versuche.\\
Aus dem Term a $\cdot 0,44^{1} \cdot 0,56^{b}$ folgt dass die Trefferwahrscheinlichkeit $p=0,44$ ist und $\mathrm{k}=1$. Daraus folgt $\mathrm{b}=\mathrm{n}-\mathrm{k}=9$ und $\mathrm{a}=\binom{10}{1}=10$.\\
Wegen $p=0,44$ werden als Treffer die Gäste gezählt, die mit Mountainbike kommen.\\
Der Term stellt die Wahrscheinlichkeit dar, dass von 10 Gästen genau ein Gast oder gar kein Gast (also höchstens ein Gast) mit dem Mountainbike kommt.

Umgekehrt könnte man es auch so formulieren, dass von 10 Gästen genau 9 oder alle 10 Gäste (also mindestens 9 Gäste) nicht mit dem Mountainbike kommen.\\
e) Die Zufallsgröße X beschreibt die Masse einer Banane in Gramm.\\
$P(0,9 \cdot 120 \leq X \leq 1,1 \cdot 120)=P(108 \leq X \leq 132) \approx 0,77$\\
f) Bedingung: $\mathrm{P}(\mathrm{X}>\mathrm{m})=0,69$\\
$\Leftrightarrow 1-P(X \leq m)=0,69$\\
$\Leftrightarrow P(X \leq m)=0,31$\\
Mit dem TR-Befehl „normInv" oder „invNorm" erhält man den Wert m $\approx 115$.\\
g) Die Zufallsgröße Y beschreibt die Masse einer Birne in Gramm.

Bedingung: $\mathrm{P}(\mathrm{Y} \leq 121) \approx 0,1$\\
Nun müssen mit dem Taschenrechner Werte für $\mu$ ausprobiert werden, sodass die obige Bedingung erfüllt ist.\\
$\mu=152: P(Y \leq 121) \approx 0,107$\\
$\mu=153: P(Y \leq 121) \approx 0,100$\\
$\mu=154: P(Y \leq 121) \approx 0,093$\\
Daraus folgt als Erwartungswert $\mu=153$.\\
Damit gilt $P(Y \geq 180) \approx 0,14$.\\
h) Die Wahrscheinlichkeit, dass die Masse eines Apfels kleiner als $\mu_{\mathrm{A}}$ ist, beträgt 0,5 (aufgrund der Symmetrie der Glockenkurve).

Die Wahrscheinlichkeit, dass die Masse eines Apfels größer als $\mu_{\mathrm{A}}$ ist, beträgt entsprechend 0,5.

Die Zufallsgröße Z zählt die Äpfel mit einer Masse, die kleiner als $\mu_{\mathrm{A}}$ sind. $Z$ ist binomialverteilt mit $n=4$ und $p=0,5$.\\
$P(Z=2)=0,375$.\\
Die Wahrscheinlichkeit, dass zwei Äpfel eine kleinere und zwei Äpfel eine größere Masse als $\mu_{\mathrm{A}}$ besitzen beträgt 0,375 .


\end{document}